% Options for packages loaded elsewhere
\PassOptionsToPackage{unicode}{hyperref}
\PassOptionsToPackage{hyphens}{url}
\PassOptionsToPackage{dvipsnames,svgnames,x11names}{xcolor}
%
\documentclass[
  letterpaper,
  DIV=11,
  numbers=noendperiod]{scrartcl}

\usepackage{amsmath,amssymb}
\usepackage{iftex}
\ifPDFTeX
  \usepackage[T1]{fontenc}
  \usepackage[utf8]{inputenc}
  \usepackage{textcomp} % provide euro and other symbols
\else % if luatex or xetex
  \usepackage{unicode-math}
  \defaultfontfeatures{Scale=MatchLowercase}
  \defaultfontfeatures[\rmfamily]{Ligatures=TeX,Scale=1}
\fi
\usepackage{lmodern}
\ifPDFTeX\else  
    % xetex/luatex font selection
\fi
% Use upquote if available, for straight quotes in verbatim environments
\IfFileExists{upquote.sty}{\usepackage{upquote}}{}
\IfFileExists{microtype.sty}{% use microtype if available
  \usepackage[]{microtype}
  \UseMicrotypeSet[protrusion]{basicmath} % disable protrusion for tt fonts
}{}
\makeatletter
\@ifundefined{KOMAClassName}{% if non-KOMA class
  \IfFileExists{parskip.sty}{%
    \usepackage{parskip}
  }{% else
    \setlength{\parindent}{0pt}
    \setlength{\parskip}{6pt plus 2pt minus 1pt}}
}{% if KOMA class
  \KOMAoptions{parskip=half}}
\makeatother
\usepackage{xcolor}
\setlength{\emergencystretch}{3em} % prevent overfull lines
\setcounter{secnumdepth}{-\maxdimen} % remove section numbering
% Make \paragraph and \subparagraph free-standing
\makeatletter
\ifx\paragraph\undefined\else
  \let\oldparagraph\paragraph
  \renewcommand{\paragraph}{
    \@ifstar
      \xxxParagraphStar
      \xxxParagraphNoStar
  }
  \newcommand{\xxxParagraphStar}[1]{\oldparagraph*{#1}\mbox{}}
  \newcommand{\xxxParagraphNoStar}[1]{\oldparagraph{#1}\mbox{}}
\fi
\ifx\subparagraph\undefined\else
  \let\oldsubparagraph\subparagraph
  \renewcommand{\subparagraph}{
    \@ifstar
      \xxxSubParagraphStar
      \xxxSubParagraphNoStar
  }
  \newcommand{\xxxSubParagraphStar}[1]{\oldsubparagraph*{#1}\mbox{}}
  \newcommand{\xxxSubParagraphNoStar}[1]{\oldsubparagraph{#1}\mbox{}}
\fi
\makeatother

\usepackage{color}
\usepackage{fancyvrb}
\newcommand{\VerbBar}{|}
\newcommand{\VERB}{\Verb[commandchars=\\\{\}]}
\DefineVerbatimEnvironment{Highlighting}{Verbatim}{commandchars=\\\{\}}
% Add ',fontsize=\small' for more characters per line
\usepackage{framed}
\definecolor{shadecolor}{RGB}{241,243,245}
\newenvironment{Shaded}{\begin{snugshade}}{\end{snugshade}}
\newcommand{\AlertTok}[1]{\textcolor[rgb]{0.68,0.00,0.00}{#1}}
\newcommand{\AnnotationTok}[1]{\textcolor[rgb]{0.37,0.37,0.37}{#1}}
\newcommand{\AttributeTok}[1]{\textcolor[rgb]{0.40,0.45,0.13}{#1}}
\newcommand{\BaseNTok}[1]{\textcolor[rgb]{0.68,0.00,0.00}{#1}}
\newcommand{\BuiltInTok}[1]{\textcolor[rgb]{0.00,0.23,0.31}{#1}}
\newcommand{\CharTok}[1]{\textcolor[rgb]{0.13,0.47,0.30}{#1}}
\newcommand{\CommentTok}[1]{\textcolor[rgb]{0.37,0.37,0.37}{#1}}
\newcommand{\CommentVarTok}[1]{\textcolor[rgb]{0.37,0.37,0.37}{\textit{#1}}}
\newcommand{\ConstantTok}[1]{\textcolor[rgb]{0.56,0.35,0.01}{#1}}
\newcommand{\ControlFlowTok}[1]{\textcolor[rgb]{0.00,0.23,0.31}{\textbf{#1}}}
\newcommand{\DataTypeTok}[1]{\textcolor[rgb]{0.68,0.00,0.00}{#1}}
\newcommand{\DecValTok}[1]{\textcolor[rgb]{0.68,0.00,0.00}{#1}}
\newcommand{\DocumentationTok}[1]{\textcolor[rgb]{0.37,0.37,0.37}{\textit{#1}}}
\newcommand{\ErrorTok}[1]{\textcolor[rgb]{0.68,0.00,0.00}{#1}}
\newcommand{\ExtensionTok}[1]{\textcolor[rgb]{0.00,0.23,0.31}{#1}}
\newcommand{\FloatTok}[1]{\textcolor[rgb]{0.68,0.00,0.00}{#1}}
\newcommand{\FunctionTok}[1]{\textcolor[rgb]{0.28,0.35,0.67}{#1}}
\newcommand{\ImportTok}[1]{\textcolor[rgb]{0.00,0.46,0.62}{#1}}
\newcommand{\InformationTok}[1]{\textcolor[rgb]{0.37,0.37,0.37}{#1}}
\newcommand{\KeywordTok}[1]{\textcolor[rgb]{0.00,0.23,0.31}{\textbf{#1}}}
\newcommand{\NormalTok}[1]{\textcolor[rgb]{0.00,0.23,0.31}{#1}}
\newcommand{\OperatorTok}[1]{\textcolor[rgb]{0.37,0.37,0.37}{#1}}
\newcommand{\OtherTok}[1]{\textcolor[rgb]{0.00,0.23,0.31}{#1}}
\newcommand{\PreprocessorTok}[1]{\textcolor[rgb]{0.68,0.00,0.00}{#1}}
\newcommand{\RegionMarkerTok}[1]{\textcolor[rgb]{0.00,0.23,0.31}{#1}}
\newcommand{\SpecialCharTok}[1]{\textcolor[rgb]{0.37,0.37,0.37}{#1}}
\newcommand{\SpecialStringTok}[1]{\textcolor[rgb]{0.13,0.47,0.30}{#1}}
\newcommand{\StringTok}[1]{\textcolor[rgb]{0.13,0.47,0.30}{#1}}
\newcommand{\VariableTok}[1]{\textcolor[rgb]{0.07,0.07,0.07}{#1}}
\newcommand{\VerbatimStringTok}[1]{\textcolor[rgb]{0.13,0.47,0.30}{#1}}
\newcommand{\WarningTok}[1]{\textcolor[rgb]{0.37,0.37,0.37}{\textit{#1}}}

\providecommand{\tightlist}{%
  \setlength{\itemsep}{0pt}\setlength{\parskip}{0pt}}\usepackage{longtable,booktabs,array}
\usepackage{calc} % for calculating minipage widths
% Correct order of tables after \paragraph or \subparagraph
\usepackage{etoolbox}
\makeatletter
\patchcmd\longtable{\par}{\if@noskipsec\mbox{}\fi\par}{}{}
\makeatother
% Allow footnotes in longtable head/foot
\IfFileExists{footnotehyper.sty}{\usepackage{footnotehyper}}{\usepackage{footnote}}
\makesavenoteenv{longtable}
\usepackage{graphicx}
\makeatletter
\newsavebox\pandoc@box
\newcommand*\pandocbounded[1]{% scales image to fit in text height/width
  \sbox\pandoc@box{#1}%
  \Gscale@div\@tempa{\textheight}{\dimexpr\ht\pandoc@box+\dp\pandoc@box\relax}%
  \Gscale@div\@tempb{\linewidth}{\wd\pandoc@box}%
  \ifdim\@tempb\p@<\@tempa\p@\let\@tempa\@tempb\fi% select the smaller of both
  \ifdim\@tempa\p@<\p@\scalebox{\@tempa}{\usebox\pandoc@box}%
  \else\usebox{\pandoc@box}%
  \fi%
}
% Set default figure placement to htbp
\def\fps@figure{htbp}
\makeatother

\KOMAoption{captions}{tableheading}
\makeatletter
\@ifpackageloaded{caption}{}{\usepackage{caption}}
\AtBeginDocument{%
\ifdefined\contentsname
  \renewcommand*\contentsname{Table of contents}
\else
  \newcommand\contentsname{Table of contents}
\fi
\ifdefined\listfigurename
  \renewcommand*\listfigurename{List of Figures}
\else
  \newcommand\listfigurename{List of Figures}
\fi
\ifdefined\listtablename
  \renewcommand*\listtablename{List of Tables}
\else
  \newcommand\listtablename{List of Tables}
\fi
\ifdefined\figurename
  \renewcommand*\figurename{Figure}
\else
  \newcommand\figurename{Figure}
\fi
\ifdefined\tablename
  \renewcommand*\tablename{Table}
\else
  \newcommand\tablename{Table}
\fi
}
\@ifpackageloaded{float}{}{\usepackage{float}}
\floatstyle{ruled}
\@ifundefined{c@chapter}{\newfloat{codelisting}{h}{lop}}{\newfloat{codelisting}{h}{lop}[chapter]}
\floatname{codelisting}{Listing}
\newcommand*\listoflistings{\listof{codelisting}{List of Listings}}
\makeatother
\makeatletter
\makeatother
\makeatletter
\@ifpackageloaded{caption}{}{\usepackage{caption}}
\@ifpackageloaded{subcaption}{}{\usepackage{subcaption}}
\makeatother

\usepackage{bookmark}

\IfFileExists{xurl.sty}{\usepackage{xurl}}{} % add URL line breaks if available
\urlstyle{same} % disable monospaced font for URLs
\hypersetup{
  pdftitle={Reproducible Research Setup},
  pdfauthor={WU Vienna \textbar{} SS 2026},
  colorlinks=true,
  linkcolor={blue},
  filecolor={Maroon},
  citecolor={Blue},
  urlcolor={Blue},
  pdfcreator={LaTeX via pandoc}}


\title{Reproducible Research Setup}
\usepackage{etoolbox}
\makeatletter
\providecommand{\subtitle}[1]{% add subtitle to \maketitle
  \apptocmd{\@title}{\par {\large #1 \par}}{}{}
}
\makeatother
\subtitle{Session 4 --- ExInt II: Research Designs in SME Research}
\author{WU Vienna \textbar{} SS 2026}
\date{2026-05-06}

\begin{document}
\maketitle


\subsection{Today's Session at a
Glance}\label{todays-session-at-a-glance}

\textbf{What we cover}

\begin{enumerate}
\def\labelenumi{\arabic{enumi}.}
\tightlist
\item
  Why reproducible research?
\item
  Project folder structure
\item
  Python environment setup (live)
\item
  Zotero \& reference management
\item
  Research question \& hypothesis
\end{enumerate}

\textbf{End-of-session deliverable}

✅ Working Python environment\\
✅ Initialized project folder on GitHub\\
✅ Zotero library with ≥ 3 references\\
✅ Draft research question + hypothesis in \texttt{README.md}

\begin{center}\rule{0.5\linewidth}{0.5pt}\end{center}

\subsection{The Replication Crisis --- Why This
Matters}\label{the-replication-crisis-why-this-matters}

\begin{quote}
\emph{``We could not reproduce 60\% of results in a sample of 100
psychology studies.''}\\
--- Open Science Collaboration (2015), \emph{Science}
\end{quote}

\textbf{The problem in business research}

\begin{itemize}
\tightlist
\item
  Undisclosed model specifications
\item
  Cherry-picked control variables
\item
  Lost data \& code after publication
\item
  No audit trail from raw data → result
\end{itemize}

\textbf{The reproducibility stack}

\begin{verbatim}
Raw Data  ──→  Code  ──→  Output
   ↑             ↑            ↑
version       version      version
controlled   controlled   controlled
\end{verbatim}

💡 \textbf{Reproducible} = someone else (or future you) can re-run your
code and get the identical result.

\begin{center}\rule{0.5\linewidth}{0.5pt}\end{center}

\subsection{What Does a Reproducible Project Look
Like?}\label{what-does-a-reproducible-project-look-like}

\begin{verbatim}
sme-internationalization/          ← Git repository
│
├── data/
│   ├── raw/          ← NEVER edit these files
│   └── processed/    ← outputs of cleaning scripts
│
├── code/
│   ├── 01_clean.py
│   ├── 02_descriptives.py
│   └── 03_regression.py
│
├── output/
│   ├── tables/
│   └── figures/
│
├── references/       ← PDFs stored here; Zotero points here
│   └── library.bib  ← auto-exported BibTeX
│
├── research_note.md  ← your main output document
├── requirements.txt  ← exact package versions
└── README.md         ← research question, hypothesis, instructions
\end{verbatim}

\textbf{Rule:} Raw data is read-only. Every transformation is a script.
Nothing lives ``in Excel''.

\begin{center}\rule{0.5\linewidth}{0.5pt}\end{center}

\subsection{Step 1 --- Setting Up Your
Environment}\label{step-1-setting-up-your-environment}

\subsubsection{Live Demo: Let's do this
together}\label{live-demo-lets-do-this-together}

\textbf{Option A: VS Code}

\begin{Shaded}
\begin{Highlighting}[]
\CommentTok{\# 1. Create project folder}
\FunctionTok{mkdir}\NormalTok{ sme{-}internationalization}
\BuiltInTok{cd}\NormalTok{ sme{-}internationalization}

\CommentTok{\# 2. Create virtual environment}
\ExtensionTok{python} \AttributeTok{{-}m}\NormalTok{ venv .venv}
\BuiltInTok{source}\NormalTok{ .venv/bin/activate   }\CommentTok{\# Mac/Linux}
\CommentTok{\# .venv\textbackslash{}Scripts\textbackslash{}activate    \# Windows}

\CommentTok{\# 3. Install packages}
\ExtensionTok{pip}\NormalTok{ install pandas numpy }\DataTypeTok{\textbackslash{}}
\NormalTok{  statsmodels linearmodels }\DataTypeTok{\textbackslash{}}
\NormalTok{  matplotlib seaborn wrds }\DataTypeTok{\textbackslash{}}
\NormalTok{  jupyter}
\end{Highlighting}
\end{Shaded}

\textbf{Option B: JupyterLab}

\begin{Shaded}
\begin{Highlighting}[]
\CommentTok{\# Same venv setup, then:}
\ExtensionTok{pip}\NormalTok{ install jupyterlab}

\CommentTok{\# Launch}
\ExtensionTok{jupyter}\NormalTok{ lab}
\end{Highlighting}
\end{Shaded}

\textbf{Key principle:} Every project gets its own virtual environment.
Never install into system Python.

\begin{center}\rule{0.5\linewidth}{0.5pt}\end{center}

\subsection{Step 2 --- Freeze Your
Dependencies}\label{step-2-freeze-your-dependencies}

\begin{Shaded}
\begin{Highlighting}[]
\CommentTok{\# After installing everything:}
\ExtensionTok{pip}\NormalTok{ freeze }\OperatorTok{\textgreater{}}\NormalTok{ requirements.txt}
\end{Highlighting}
\end{Shaded}

\begin{Shaded}
\begin{Highlighting}[]
\NormalTok{\# requirements.txt (example)}
\NormalTok{pandas==2.2.1}
\NormalTok{numpy==1.26.4}
\NormalTok{statsmodels==0.14.1}
\NormalTok{linearmodels==5.4}
\NormalTok{matplotlib==3.8.3}
\NormalTok{seaborn==0.13.2}
\NormalTok{wrds==3.2.0}
\end{Highlighting}
\end{Shaded}

Anyone can now reproduce your environment exactly:

\begin{Shaded}
\begin{Highlighting}[]
\ExtensionTok{pip}\NormalTok{ install }\AttributeTok{{-}r}\NormalTok{ requirements.txt}
\end{Highlighting}
\end{Shaded}

\textbf{Why does this matter?} Package updates break code.
\texttt{statsmodels} 0.13 → 0.14 changed default standard errors.
Without pinned versions, your results may not reproduce even on your own
machine next year.

\begin{center}\rule{0.5\linewidth}{0.5pt}\end{center}

\subsection{Step 3 --- Git \& GitHub}\label{step-3-git-github}

\begin{Shaded}
\begin{Highlighting}[]
\CommentTok{\# Initialize repository}
\FunctionTok{git}\NormalTok{ init}
\FunctionTok{git}\NormalTok{ add .}
\FunctionTok{git}\NormalTok{ commit }\AttributeTok{{-}m} \StringTok{"initial project structure"}

\CommentTok{\# Push to GitHub}
\ExtensionTok{gh}\NormalTok{ repo create sme{-}internationalization }\AttributeTok{{-}{-}public}
\FunctionTok{git}\NormalTok{ push }\AttributeTok{{-}u}\NormalTok{ origin main}
\end{Highlighting}
\end{Shaded}

\textbf{What goes in \texttt{.gitignore}?}

\begin{Shaded}
\begin{Highlighting}[]
\NormalTok{\# .gitignore}
\NormalTok{.venv/}
\NormalTok{data/raw/          \# raw data often has licensing restrictions}
\NormalTok{*.pyc}
\NormalTok{\_\_pycache\_\_/}
\NormalTok{.DS\_Store}
\NormalTok{*.egg{-}info/}
\end{Highlighting}
\end{Shaded}

📌 \textbf{Your WRDS credentials never go in your code.} Use environment
variables or a \texttt{.env} file (also in \texttt{.gitignore}).

\begin{center}\rule{0.5\linewidth}{0.5pt}\end{center}

\subsection{Step 4 --- Reference Management with
Zotero}\label{step-4-reference-management-with-zotero}

\textbf{Setup (5 minutes)}

\begin{enumerate}
\def\labelenumi{\arabic{enumi}.}
\tightlist
\item
  Download Zotero (free) + browser connector
\item
  Set storage folder → \texttt{references/} in your project
\item
  Install \textbf{Better BibTeX} plugin
\item
  Configure auto-export to \texttt{references/library.bib}
\end{enumerate}

\textbf{Workflow}

\begin{itemize}
\tightlist
\item
  Find paper → click browser connector → PDF saved automatically
\item
  Zotero manages metadata
\item
  \texttt{library.bib} is always up to date
\item
  Cite in Markdown: \texttt{{[}@Johanson1977{]}}
\end{itemize}

\textbf{In your \texttt{research\_note.md}:}

\begin{Shaded}
\begin{Highlighting}[]
\CommentTok{{-}{-}{-}}
\AnnotationTok{bibliography:}\CommentTok{ references/library.bib}
\CommentTok{{-}{-}{-}}

\NormalTok{SMEs face liability of foreignness}
\CommentTok{[}\OtherTok{@Zaheer1995}\CommentTok{]}\NormalTok{, which increases...}

\NormalTok{The Uppsala model }\CommentTok{[}\OtherTok{@Johanson1977}\CommentTok{]}
\NormalTok{suggests that firms...}
\end{Highlighting}
\end{Shaded}

Render with Quarto or Pandoc:

\begin{Shaded}
\begin{Highlighting}[]
\ExtensionTok{quarto}\NormalTok{ render research\_note.md}
\CommentTok{\# → research\_note.pdf / .html}
\end{Highlighting}
\end{Shaded}

\begin{center}\rule{0.5\linewidth}{0.5pt}\end{center}

\subsection{Step 5 --- The Research
Question}\label{step-5-the-research-question}

\textbf{A good research question is\ldots{}}

\begin{itemize}
\tightlist
\item
  \textbf{Focused}: one clear phenomenon
\item
  \textbf{Grounded}: linked to existing theory
\item
  \textbf{Falsifiable}: can be answered with data
\item
  \textbf{Relevant}: contributes to a gap
\end{itemize}

\textbf{Template:}

\begin{quote}
\emph{``How does {[}X{]} affect {[}Y{]} among {[}population{]}, and does
{[}Z{]} moderate this relationship?''}
\end{quote}

\textbf{Example for our project:}

\emph{``Does degree of internationalization (DOI) exhibit a non-linear
(inverted U-shaped) relationship with firm performance among European
SMEs, and does R\&D intensity moderate this relationship?''}

\textbf{Theoretical grounding:} - Internalization theory (Buckley \&
Casson 1976) - Liability of foreignness (Zaheer 1995) -
Learning-by-exporting (Clerides et al.~1998)

\begin{center}\rule{0.5\linewidth}{0.5pt}\end{center}

\subsection{Step 6 --- From Theory to
Hypothesis}\label{step-6-from-theory-to-hypothesis}

\textbf{The logic chain:}

\begin{verbatim}
Theory  →  Theoretical argument  →  Testable prediction  →  Falsifiable hypothesis
\end{verbatim}

\textbf{Theoretical argument:}

\begin{quote}
Initial internationalization generates learning benefits and market
diversification gains. Beyond a threshold, coordination costs, cultural
distance costs, and managerial complexity outweigh these benefits.
\end{quote}

→ Inverted U-shape prediction

\textbf{Formal hypothesis:}

\textbf{H1:} DOI has a positive effect on SME performance at low levels
of internationalization, but this effect turns negative at high levels
(inverted U-shaped relationship).

\textbf{H2:} R\&D-intensive SMEs sustain performance gains at higher DOI
levels than low R\&D-intensity SMEs.

\begin{center}\rule{0.5\linewidth}{0.5pt}\end{center}

\subsection{Your README.md --- What It Must
Contain}\label{your-readme.md-what-it-must-contain}

\begin{Shaded}
\begin{Highlighting}[]
\FunctionTok{\# SME Internationalization \& Firm Performance}
\FunctionTok{\#\#\# ExInt II Research Project | WU Vienna | SS 2026}
\NormalTok{**Author:** }\CommentTok{[}\OtherTok{Your Name}\CommentTok{]}

\FunctionTok{\#\# Research Question}
\NormalTok{Does degree of internationalization exhibit an inverted U{-}shaped relationship}
\NormalTok{with firm performance among European SMEs, and does R\&D intensity moderate}
\NormalTok{this relationship?}

\FunctionTok{\#\# Hypotheses}
\SpecialStringTok{{-} }\NormalTok{**H1:** DOI has a positive effect on performance at low levels but turns}
\NormalTok{  negative at high levels (inverted U{-}shape).}
\SpecialStringTok{{-} }\NormalTok{**H2:** R\&D intensity positively moderates the DOI–performance relationship.}

\FunctionTok{\#\# Data}
\SpecialStringTok{{-} }\NormalTok{Source: WRDS / Compustat Global}
\SpecialStringTok{{-} }\NormalTok{Sample: European SMEs (employees \textless{} 250 or assets \textless{} €43m), 2005–2020}
\SpecialStringTok{{-} }\NormalTok{Key variables: }\InformationTok{\textasciigrave{}at\textasciigrave{}}\NormalTok{ (total assets), }\InformationTok{\textasciigrave{}sale\textasciigrave{}}\NormalTok{ (sales), }\InformationTok{\textasciigrave{}xrd\textasciigrave{}}\NormalTok{ (R\&D expenditure)}

\FunctionTok{\#\# How to Reproduce}
\SpecialStringTok{1. }\InformationTok{\textasciigrave{}git clone https://github.com/[you]/sme{-}internationalization\textasciigrave{}}
\SpecialStringTok{2. }\InformationTok{\textasciigrave{}pip install {-}r requirements.txt\textasciigrave{}}
\SpecialStringTok{3. }\NormalTok{Add WRDS credentials to }\InformationTok{\textasciigrave{}.env\textasciigrave{}}
\SpecialStringTok{4. }\NormalTok{Run scripts in order: }\InformationTok{\textasciigrave{}01\_clean.py\textasciigrave{}}\NormalTok{ → }\InformationTok{\textasciigrave{}02\_descriptives.py\textasciigrave{}}\NormalTok{ → }\InformationTok{\textasciigrave{}03\_regression.py\textasciigrave{}}
\end{Highlighting}
\end{Shaded}

\begin{center}\rule{0.5\linewidth}{0.5pt}\end{center}

\subsection{Session 4 Checklist}\label{session-4-checklist}

Before you leave today, confirm:

\begin{longtable}[]{@{}
  >{\raggedright\arraybackslash}p{(\linewidth - 2\tabcolsep) * \real{0.4615}}
  >{\raggedright\arraybackslash}p{(\linewidth - 2\tabcolsep) * \real{0.5385}}@{}}
\toprule\noalign{}
\begin{minipage}[b]{\linewidth}\raggedright
Task
\end{minipage} & \begin{minipage}[b]{\linewidth}\raggedright
Done?
\end{minipage} \\
\midrule\noalign{}
\endhead
\bottomrule\noalign{}
\endlastfoot
Virtual environment created \& activated & ☐ \\
All packages installed
(\texttt{pip\ freeze\ \textgreater{}\ requirements.txt}) & ☐ \\
Folder structure created (\texttt{data/}, \texttt{code/},
\texttt{output/}, \texttt{references/}) & ☐ \\
Git repo initialized \& pushed to GitHub & ☐ \\
\texttt{.gitignore} in place & ☐ \\
Zotero connected, Better BibTeX installed & ☐ \\
≥ 3 papers in Zotero, \texttt{library.bib} auto-exporting to
\texttt{references/} & ☐ \\
Research question drafted in \texttt{README.md} & ☐ \\
At least one hypothesis written in \texttt{README.md} & ☐ \\
\end{longtable}

\begin{center}\rule{0.5\linewidth}{0.5pt}\end{center}

\subsection{Looking Ahead --- Sessions
5--7}\label{looking-ahead-sessions-57}

\begin{longtable}[]{@{}
  >{\raggedright\arraybackslash}p{(\linewidth - 2\tabcolsep) * \real{0.5625}}
  >{\raggedright\arraybackslash}p{(\linewidth - 2\tabcolsep) * \real{0.4375}}@{}}
\toprule\noalign{}
\begin{minipage}[b]{\linewidth}\raggedright
Session
\end{minipage} & \begin{minipage}[b]{\linewidth}\raggedright
Focus
\end{minipage} \\
\midrule\noalign{}
\endhead
\bottomrule\noalign{}
\endlastfoot
\textbf{Session 5} & WRDS API access · Data pull · Cleaning pipeline \\
\textbf{Session 6} & Descriptive statistics · Panel structure · Fixed
effects \\
\textbf{Session 7} & Regression results · Robustness checks ·
Interpretation \\
\textbf{Session 8} & Research note writing in Markdown · Peer review ·
Submission \\
\end{longtable}

\textbf{Your final deliverable} is \texttt{research\_note.md} --- a
complete, self-contained research note rendered to PDF via Quarto. Raw
data → code → output → written results, fully traceable.

\begin{center}\rule{0.5\linewidth}{0.5pt}\end{center}

\subsection{Resources}\label{resources}

\textbf{Tools}

\begin{itemize}
\tightlist
\item
  \href{https://www.zotero.org/}{Zotero} +
  \href{https://retorque.re/zotero-better-bibtex/}{Better BibTeX}
\item
  \href{https://quarto.org/}{Quarto} for rendering Markdown to PDF/HTML
\item
  \href{https://code.visualstudio.com/}{VS Code} + Python extension
\item
  \href{https://wrds-www.wharton.upenn.edu/pages/support/programming-wrds/programming-python/}{WRDS
  Python package docs}
\end{itemize}

\textbf{Key papers for your \texttt{library.bib}}

\begin{itemize}
\tightlist
\item
  Johanson \& Vahlne (1977) --- Uppsala model
\item
  Zaheer (1995) --- Liability of foreignness\\
\item
  Lu \& Beamish (2001) --- SME internationalization \& performance
\item
  Hitt et al.~(1997) --- International diversification \& performance
\end{itemize}

📁 Example project available at:
\texttt{github.com/{[}instructor{]}/exint2-example-project}




\end{document}
